\section{Введение}
\label{sec:heldout-context}
Система для разработки программ должна интерпретировать спецификацию, использовать относящиеся к задаче код и документацию, построить решение и проверить его тестами. Эти функции можно распределить между этапами планирования, реализации и проверки. Если отдельные вызовы модели получают одни и те же материалы, контекст приходится передавать повторно. Мы исследуем, можно ли сохранить эти функции при меньшем числе вызовов, оставив проверку корректности внешним тестам.

Hybrid-INoT отвечает на этот вопрос сочетанием нескольких предписанных функций в одном ответе модели и последующей проверки программы исполняемыми тестами. Планирование, реализация и проверка задаются в общем вызове модели, а внешняя процедура оценивает полученную программу. Дизайн развивает мотивацию внутреннего диалога INoT~\cite{sun2025inot}, но исследуемая ролевая конфигурация отличается от исходного алгоритма PromptCode. Она отделяет генерацию программы от проверки её корректности. Повторная передача длинного общего контекста мотивирует анализ затрат, но сама по себе не гарантирует сохранения качества при одном вызове.

Исходная реализация Hybrid-INoT включала также отбор контекста и избирательные повторы. Сравнение целых конфигураций изменяло эти механизмы вместе с топологией вызовов и потому не выделяет вклад внутренних ролей. Настоящая оценка изолирует одношаговое ядро архитектуры: полный предоставленный контекст, один итоговый кандидат, предписанные планирование, реализацию и проверку, а также последующее штатное тестирование. Компрессия и повторы по результатам проверки отключены. Раздельная оценка компонентов позволяет проверить проектные решения Hybrid-INoT в рамках исходной инженерной задачи.

У ядра есть два разных эмпирических вопроса. Меняют ли названия ролей результат при неизменных операциях этапов? Как выполнение этих операций в одном ответе соотносится по качеству и ресурсам с их распределением по трём вызовам? Прямой решатель проверяет полезность поэтапной организации в целом. Отдельная адаптация INoT исследует родственную алгоритмическую альтернативу. Каждое условие получает одинаковые задачу и контекст и оценивается одной последующей процедурой. Внутренняя проверка, описанная в запросе, не считается отдельным измеренным исходом.

Вклад работы включает:
\begin{enumerate}
\item Операционное описание однопроходного ядра Hybrid-INoT, его состояния и границы внешней оценки с соответствием реализованным контрольным условиям.
\item Условный баланс токенов общего контекста, который объясняет возможную экономию при меньшем числе вызовов без предположения о равной длине ответов.
\item Факторный эксперимент $2\times2$ без компрессии с прямым решателем на 1\,000 задачах BigCodeBench, тремя повторами каждого условия и 14\,961 программой, проверенной штатными тестами. Четыре заранее заданных сравнения отделяют ролевые обозначения от структуры вызовов.
\item Дополнительные данные первого факторного эксперимента на 200 задачах, исследования обозначений и разметки на 80 задачах, разработки, адаптации INoT и демонстрации исправления одного репозитория. В приложениях сохранены инструкции и все повторные исходы.
\end{enumerate}
Абляции ролей и исследования числа взаимодействий имеют близких предшественников; вклад работы заключается в определённой архитектурной оценке и её проверяемых данных. Большой эксперимент показывает увеличение расхода и снижение успеха тестов при разделении этапов на три вызова. Общая польза названий ролей для качества и неухудшение качества не установлены.

\section{Связанные исследования}
\label{sec:heldout-related}
\subsection{Ближайший методологический предшественник}
В работе Self-collaboration Code Generation via ChatGPT~\cite{dong2023selfcollaboration} уже исследованы ролевые инструкции и максимум взаимодействий в разделах~5.2--5.3. Поэтому удаление ролей или изменение взаимодействий не заявляется новизной настоящего исследования. Таблица~\ref{tab:closest-method} разграничивает реализованные воздействия.

\begin{table}[htbp]\centering\small
\caption{Сопоставление с ближайшим предшественником~\cite{dong2023selfcollaboration}, таблицы~3--4 и приложение~B.1. Раунды взаимодействия и новые вызовы модели являются разными величинами.}\label{tab:closest-method}
\begin{tabular}{p{0.20\textwidth}p{0.31\textwidth}p{0.36\textwidth}}\toprule
Аспект & Self-collaboration & Настоящий эксперимент\\\midrule
Обозначения & Ролевые инструкции и инструкции без ролей & Одинаковые предложения об операциях; меняются только префиксы этапов\\
Взаимодействие & Цикл обратной связи; варьируется максимум взаимодействий & Один или три новых вызова; фиксированная последовательность, полная видимая история\\
Дизайн & Раздельные абляции ролей и взаимодействия & Совместная матрица обозначений и вызовов для каждой отобранной задачи\\
Задачи & HumanEval, MBPP и расширенные тесты & Резервная выборка BigCodeBench; эталонные и неправильные контроли до генерации\\\bottomrule
\end{tabular}\end{table}

Добавленная ценность состоит в репликации и расширении с акцентом на проверяемость: совместный дизайн оценивает взаимодействие обозначений и топологии при одном интерфейсе генерации, а связь результатов штатных тестов с проверенными программами и компоненты токенов делают различия качества и ресурсов проверяемыми. Это новые эмпирические условия и набор доказательств, а не новый принцип ролевого сотрудничества или первая абляция ролей.

Для внешнего сравнения наиболее близок к вопросу об обозначениях метод Self-collaboration вместе с вариантом без ролей. Его реализация должна сохранять обратную связь и правило остановки; переименование нашего фиксированного трёхвызовного MR не воспроизводит этот метод. Отдельный пилот авторской реализации на трёх задачах описан в разделе~\ref{sec:scc-feasibility}; внешнее сравнение с достаточной мощностью и превосходство над SCC остаются непроверенными. В этой схеме D задаёт минимальное реализованное сравнение, а SN и MN служат контролями для выделения эффекта обозначений. INoT* отвечает на отдельный вопрос об алгоритме внутреннего диалога и поэтому не рассматривается как ближайший предшественник абляции ролей.

\subsection{Внутренний диалог и репозиторные методы}

INoT описывает компактные инструкции PromptCode для внутреннего диалога и рассуждения~\cite{sun2025inot}. Это мотивирует оценку раскрытой адаптации алгоритма, тогда как отдельное минимальное воздействие обозначений проверяет, меняют ли результат сами названия ролей. Сравнения отвечают на разные вопросы: инструкция INoT изменяет процедуру рассуждения, а факторный контраст обозначений сохраняет предложения с описанием операций. Исходное сравнение Hybrid-INoT дополнительно изменяло компрессию и повторные попытки; его прежние условия рассматриваются только как контекст и не объединяются с новым экспериментом.

Исследование одно- и многоагентных систем при сопоставимом бюджете токенов рассуждения подчёркивает различие дополнительных вызовов и вычислений~\cite{tran2026budget}. Поскольку здесь длина завершения не ограничена общим жёстким бюджетом, токены и денежная оценка рассматриваются вместе с качеством как исходы реализованных протоколов; оптимальное распределение фиксированного бюджета инференса из этого дизайна не следует.

Agentless использует структурированный процесс локализации, исправления и проверки~\cite{xia2024agentless}; SWE-agent связывает модель с интерфейсом репозитория~\cite{yang2024sweagent}. Обе системы добавляют поиск, инструменты и взаимодействие с окружением сверх генерации функций при фиксированном контексте. Их опубликованные оценки нельзя вставлять в парную матрицу с другими задачами и настройками. BigCodeBench~\cite{zhuo2024bigcodebench} предоставляет разнообразные задачи с библиотечными зависимостями и исполнимыми тестами, а SWE-bench~\cite{jimenez2023swebench,swebenchdocs2026} проверяет изменения репозитория тестами конкретных задач из трекера ошибок. Штатная оценка связывает исходы с исполняемыми тестами, однако успешный эталонный контроль не разрешает каждое расхождение естественно-языковой инструкции и теста.

Обратная связь является отдельным выбором в дизайне метода. Self-Refine чередует замечания и исправления, полученные от той же модели~\cite{madaan2023selfrefine}. Reflexion сохраняет словесные выводы между попытками; в его экспериментах с программированием обратную связь дают исполняемые тесты, созданные моделью~\cite{shinn2023reflexion}. Условия нашего компонентного эксперимента не получают результатов исполнения во время генерации и следуют фиксированной последовательности, поэтому не воспроизводят эти адаптивные процедуры. Авторская реализация SCC использует сгенерированные тесты как часть своей полной процедуры. В наших экспериментах тесты бенчмарка и эталонные программы недоступны интерфейсу генерации; внутреннее исполнение SCC не определяет успех на бенчмарке. Поэтому доступ к обратной связи и остановка входят в сравнение внешних методов, а не доказывают эффект ролевых обозначений. Self-Refine и Reflexion включены в методологическое сопоставление; экспериментальные запуски этих методов здесь не выполнены.

Исследования persona prompting мотивируют более узкое утверждение, чем универсальная польза ролей для рассуждения. Zheng и соавторы не обнаруживают устойчивого улучшения от системных персон в исследованных фактологических задачах~\cite{zheng2024personas}. Luz de Araujo и соавторы разделяют качество, устойчивость к несущественным атрибутам персоны и соответствие заданной специализации; результаты зависят от модели и задачи~\cite{araujo2025personas}. Эти работы не проверяют наши запросы на генерацию кода. Они обосновывают изменение обозначений при сохранении требуемых операций и ограничение вывода конкретными использованными обозначениями.

Данные о нескольких агентах также зависят от выбранного сравнения. Choi и соавторы связывают значительную часть выигрыша в своих NLP-экспериментах с независимым голосованием, отделяя обмен сообщениями от агрегации~\cite{choi2025debate}. В то же время Wunderlich и соавторы находят конфигурации, в которых debate или mixture-of-agents улучшают соотношение вычислений и точности относительно self-consistency на MMLU-Pro и BBH~\cite{wunderlich2026compute}. Это основание для условной оценки, а не универсального вывода за или против нескольких агентов. Наша нейтральная последовательная цепочка не является голосованием, смесью независимых моделей или ансамблем с выровненным бюджетом. В работе выполнена факторная проверка одной модели на исполняемом коде с сохранением трасс и учётом доступных счётчиков ресурсов; общее превосходство над перечисленными альтернативами не входит в оцениваемый эффект.

